% =====================================================================
%  TD Séries numériques — exercices complémentaires (absents du livret)
%  Inspirés des exercices classiques du chapitre ; calculs vérifiés numériquement.
% =====================================================================

\begin{exo}{sn-1}{1}{suites classiques}
Déterminer le terme général des suites définies par :
\begin{enumerate}
\item $u_0=0$ et $u_{n+1}=-u_n+2$ ;
\item $u_0=1$, $u_1=2$ et $u_{n+2}+u_{n+1}+u_n=0$ ;
\item $u_0=0$, $u_1=1$ et $u_{n+2}-3u_{n+1}+2u_n=0$.
\end{enumerate}
\end{exo}

\begin{corrige}{sn-1}
\q{1.} Point fixe : $c=-c+2$, $c=1$. La suite $(u_n-1)$ est géométrique de raison $-1$ : $u_n-1=(-1)^n(u_0-1)$, donc
$u_n=1-(-1)^n$ (la suite vaut $0,2,0,2,\dots$).
\q{2.} $r^2+r+1=0$ a pour racines $e^{\pm2i\pi/3}$ ($\rho=1$, $\theta=\frac{2\pi}3$) : $u_n=A\cos\frac{2n\pi}3+B\sin\frac{2n\pi}3$.
$u_0=1$ donne $A=1$ ; $u_1=-\frac12+\frac{\sqrt3}2B=2$ donne $B=\frac5{\sqrt3}$. Donc $u_n=\cos\frac{2n\pi}3+\frac{5}{\sqrt3}\sin\frac{2n\pi}3$.
Contrôle : $u_2=-u_1-u_0=-3$ et la formule donne $-\frac12-\frac52=-3$.
\q{3.} $r^2-3r+2=(r-1)(r-2)$ : $u_n=A+B\,2^n$, avec $A+B=0$ et $A+2B=1$, donc $u_n=2^n-1$.
\end{corrige}

\begin{exo}{sn-2}{1}{équivalents}
Donner un équivalent simple de :
\[\text{(a)}\ \sqrt{n^2+n+1}-n\qquad\text{(b)}\ \sin\frac{n\pi}{n+1}\qquad\text{(c)}\ \sqrt{n+(-1)^n}-\sqrt n\qquad\text{(d)}\ \ln(n+1)-\ln n.\]
\end{exo}

\begin{corrige}{sn-2}
\q{(a)} $\sqrt{n^2+n+1}-n=\frac{n+1}{\sqrt{n^2+n+1}+n}\sim\frac n{2n}=\frac12$ : la suite tend vers $\frac12$ (équivalent : la constante $\frac12$).
\q{(b)} $\sin\frac{n\pi}{n+1}=\sin\left(\pi-\frac\pi{n+1}\right)=\sin\frac\pi{n+1}\sim\frac\pi n$.
\q{(c)} $\sqrt n\left(\sqrt{1+\frac{(-1)^n}n}-1\right)=\sqrt n\left(\frac{(-1)^n}{2n}+O\left(\frac1{n^2}\right)\right)\sim\frac{(-1)^n}{2\sqrt n}$.
\q{(d)} $\ln\left(1+\frac1n\right)\sim\frac1n$.
\piege{en (a), « $\sqrt{n^2+n+1}\sim n$ » ne permet pas de conclure : on ne soustrait pas des équivalents. On multiplie par la quantité conjuguée.}
\end{corrige}

\begin{exo}{sn-3}{2}{théorème des gendarmes}
Étudier la limite de $u_n=\sum_{k=1}^n\frac1{\sqrt{n^2+k}}$, puis, pour $x\in\R$ fixé, de $v_n=\frac1{n^2}\sum_{k=1}^n\lfloor kx\rfloor$.
\end{exo}

\begin{corrige}{sn-3}
\q{1.} Pour $1\leq k\leq n$ : $\frac1{\sqrt{n^2+n}}\leq\frac1{\sqrt{n^2+k}}\leq\frac1{\sqrt{n^2+1}}$. On somme les $n$ termes :
$\frac n{\sqrt{n^2+n}}\leq u_n\leq\frac n{\sqrt{n^2+1}}$, et les deux bornes tendent vers $1$ : $u_n\to1$.
\q{2.} $kx-1<\lfloor kx\rfloor\leq kx$. On somme : $x\frac{n(n+1)}2-n<\sum_{k=1}^n\lfloor kx\rfloor\leq x\frac{n(n+1)}2$, puis on divise par $n^2$ :
les deux bornes tendent vers $\frac x2$, donc $v_n\to\frac x2$.
\end{corrige}

\begin{exo}{sn-4}{2}{nature : comparaison, équivalents, d'Alembert}
Déterminer la nature des séries de terme général :
\[\text{(a)}\ \frac{2^nn!}{n^n}\quad\text{(b)}\ \frac{n^4}{3^n}\quad\text{(c)}\ \sin\frac1n+\ln\left(1-\frac1n\right)\quad\text{(d)}\ n\sin\frac1n\quad
\text{(e)}\ \sin(e^{-3n})\quad\text{(f)}\ \frac1{\sqrt[3]{n^2-1}}-\frac1{\sqrt[3]{n^2+1}}\quad\text{(g)}\ \frac1{\sqrt n\cos^2n}.\]
\end{exo}

\begin{corrige}{sn-4}
\q{(a)} $\frac{u_{n+1}}{u_n}=2\left(\frac n{n+1}\right)^n=2\left(1+\frac1n\right)^{-n}\to\frac2e<1$ : convergente (d'Alembert).
\q{(b)} $\frac{u_{n+1}}{u_n}=\frac13\left(1+\frac1n\right)^4\to\frac13$ : convergente.
\q{(c)} $\sin x=x+O(x^3)$ et $\ln(1-x)=-x-\frac{x^2}2+O(x^3)$ : $u_n=-\frac1{2n^2}+O\left(\frac1{n^3}\right)\sim-\frac1{2n^2}$. Le terme est
\emph{négatif} à partir d'un rang ; $-u_n\sim\frac1{2n^2}\geq0$ : convergente.
\q{(d)} $n\sin\frac1n\to1\neq0$ : divergence grossière.
\q{(e)} $\sin(e^{-3n})\sim e^{-3n}$, terme d'une série géométrique de raison $e^{-3}<1$ : convergente.
\q{(f)} $n^{-2/3}\left[\left(1-\frac1{n^2}\right)^{-1/3}-\left(1+\frac1{n^2}\right)^{-1/3}\right]=n^{-2/3}\left(\frac2{3n^2}+O\left(\frac1{n^4}\right)\right)\sim\frac2{3n^{8/3}}$ :
convergente ($\frac83>1$).
\q{(g)} $0<\cos^2n\leq1$ (et $\cos n\neq0$ car $n$ n'est pas de la forme $\frac\pi2+k\pi$), donc $u_n\geq\frac1{\sqrt n}$ : divergente par
comparaison.
\retenir{une série de signe constant \emph{négatif} se traite comme une série à termes positifs : on étudie $\sum(-u_n)$.}
\end{corrige}

\begin{exo}{sn-5}{1}{décomposition en éléments simples}
Justifier la convergence et calculer $\sum_{n=1}^{+\infty}\frac1{n(n+2)}$.
\end{exo}

\begin{corrige}{sn-5}
$\frac1{n(n+2)}=\frac12\left(\frac1n-\frac1{n+2}\right)$. Les termes se simplifient de deux en deux :
$\sum_{k=1}^n\frac1{k(k+2)}=\frac12\left(1+\frac12-\frac1{n+1}-\frac1{n+2}\right)\to\frac34$. La série converge et sa somme vaut $\frac34$.
\gain{écrire les premiers et les derniers termes pour voir ce qui reste : ici $\frac11$, $\frac12$ d'un côté, $\frac1{n+1}$, $\frac1{n+2}$ de l'autre.}
\end{corrige}

\begin{exo}{sn-6}{2}{série géométrique dérivée}
\begin{enumerate}
\item Pour $x\neq1$, calculer $\sum_{k=0}^nx^k$, puis $\sum_{k=1}^nkx^{k-1}$ en dérivant.
\item En déduire que, pour $|x|<1$, $\sum_{k=1}^{+\infty}kx^{k-1}=\frac1{(1-x)^2}$.
\item Justifier la convergence et calculer $\sum_{k=0}^{+\infty}\big(k-(-1)^k\big)3^{-k}$.
\end{enumerate}
\end{exo}

\begin{corrige}{sn-6}
\q{1.} $\sum_{k=0}^nx^k=\frac{1-x^{n+1}}{1-x}$. On dérive les deux membres :
$\sum_{k=1}^nkx^{k-1}=\frac{-(n+1)x^n(1-x)+(1-x^{n+1})}{(1-x)^2}=\frac{1-(n+1)x^n+nx^{n+1}}{(1-x)^2}$.
\q{2.} Si $|x|<1$, $nx^n\to0$ et $(n+1)x^n\to0$ (croissances comparées), d'où la limite $\frac1{(1-x)^2}$.
\q{3.} $\sum k3^{-k}=\frac13\sum k\left(\frac13\right)^{k-1}=\frac13\cdot\frac1{(2/3)^2}=\frac34$ et $\sum(-1)^k3^{-k}=\frac1{1+\frac13}=\frac34$. Les
deux séries convergent, donc la série étudiée aussi (linéarité), et sa somme vaut $\frac34-\frac34=0$.
\end{corrige}

\begin{exo}{sn-7}{3}{somme d'une série alternée par une intégrale}
\begin{enumerate}
\item Montrer que $\sum\frac{(-1)^n}{3n+1}$ converge.
\item Montrer que, pour tout $n$, $\displaystyle\sum_{k=0}^n\frac{(-1)^k}{3k+1}=\int_0^1\frac{dt}{1+t^3}-\int_0^1\frac{(-t^3)^{n+1}}{1+t^3}\,dt$.
\item En déduire $\displaystyle\sum_{n=0}^{+\infty}\frac{(-1)^n}{3n+1}=\int_0^1\frac{dt}{1+t^3}$, puis calculer cette intégrale
(décomposer $\frac1{1+t^3}=\frac\alpha{1+t}+\frac{\beta t+\gamma}{t^2-t+1}$).
\end{enumerate}
\end{exo}

\begin{corrige}{sn-7}
\q{1.} $\frac1{3n+1}$ est positive, décroissante, de limite nulle : critère des séries alternées.
\q{2.} $\sum_{k=0}^n(-t^3)^k=\frac{1-(-t^3)^{n+1}}{1+t^3}$ ; on intègre sur $[0,1]$, avec $\int_0^1(-1)^kt^{3k}\,dt=\frac{(-1)^k}{3k+1}$.
\q{3.} $\left|\int_0^1\frac{(-t^3)^{n+1}}{1+t^3}\,dt\right|\leq\int_0^1t^{3n+3}\,dt=\frac1{3n+4}\to0$. Ensuite $\alpha=\frac13$, $\beta=-\frac13$,
$\gamma=\frac23$ : $\frac{-t+2}{t^2-t+1}=-\frac12\cdot\frac{2t-1}{t^2-t+1}+\frac32\cdot\frac1{(t-\frac12)^2+\frac34}$. D'où
\[\int_0^1\frac{dt}{1+t^3}=\frac13\ln2-\frac16\big[\ln(t^2-t+1)\big]_0^1+\frac12\cdot\frac2{\sqrt3}\left[\arctan\frac{2t-1}{\sqrt3}\right]_0^1=\frac{\ln2}3+\frac\pi{3\sqrt3}.\]
(Numériquement $0{,}8356\dots$, vérifié.)
\end{corrige}

\begin{exo}{sn-8}{2}{séries de Bertrand}
\begin{enumerate}
\item Par comparaison avec une intégrale, montrer que $\sum_{n\geq2}\frac1{n(\ln n)^2}$ converge.
\item Montrer que $\sum_{k=2}^n\frac1{k\ln k}\sim\ln(\ln n)$.
\end{enumerate}
\end{exo}

\begin{corrige}{sn-8}
\q{1.} $f(t)=\frac1{t(\ln t)^2}$ est continue, positive, décroissante sur $[2,+\infty[$ (produit de fonctions positives croissantes au
dénominateur), et $\int_2^X f=\left[-\frac1{\ln t}\right]_2^X=\frac1{\ln2}-\frac1{\ln X}\to\frac1{\ln2}$ : l'intégrale converge, donc la série aussi.
\q{2.} $g(t)=\frac1{t\ln t}$ est décroissante sur $[2,+\infty[$ et $\int_2^ng=\ln\ln n-\ln\ln2$. L'encadrement du théorème série-intégrale
donne $\int_2^{n+1}g\leq\sum_{k=2}^ng(k)\leq g(2)+\int_2^ng$, soit $\ln\ln(n+1)-\ln\ln2\leq\sum_{k=2}^n\frac1{k\ln k}\leq\frac1{2\ln2}+\ln\ln n-\ln\ln2$.
Les deux bornes sont équivalentes à $\ln\ln n$ (car $\ln(n+1)\sim\ln n$, donc $\ln\ln(n+1)-\ln\ln n\to0$).
\end{corrige}

\begin{exo}{sn-9}{2}{une série semi-convergente}
\begin{enumerate}
\item Montrer que $\sum_{k=1}^n\frac{\ln k}k\sim\frac{(\ln n)^2}2$.
\item Montrer que $\sum\frac{(-1)^n\ln n}n$ converge, mais pas absolument.
\end{enumerate}
\end{exo}

\begin{corrige}{sn-9}
\q{1.} $f(t)=\frac{\ln t}t$ a pour dérivée $\frac{1-\ln t}{t^2}\leq0$ sur $[e,+\infty[$ : $f$ décroît sur $[3,+\infty[$. Avec $\int_3^nf=\frac{(\ln n)^2-(\ln3)^2}2$,
l'encadrement série-intégrale donne $\sum_{k=3}^n\frac{\ln k}k=\frac{(\ln n)^2}2+O(1)$ (les deux bornes diffèrent de quantités bornées), d'où
l'équivalent ; les deux premiers termes ne changent rien.
\q{2.} $a_n=\frac{\ln n}n\geq0$ tend vers $0$ et décroît pour $n\geq3$ : la série converge (critère, à partir du rang 3). Elle ne converge pas
absolument, car $\sum_{k\leq n}\frac{\ln k}k\to+\infty$ d'après 1.
\end{corrige}

\begin{exo}{sn-16}{2}{une famille de séries}
Selon $\alpha\in\R$, déterminer la nature de $\sum_{n\geq2}\frac1{n^\alpha\ln n}$.
\end{exo}

\begin{corrige}{sn-16}
\q{$\alpha>1$.} On choisit $\beta$ avec $1<\beta<\alpha$ : $n^\beta\cdot\frac1{n^\alpha\ln n}=\frac1{n^{\alpha-\beta}\ln n}\to0$, donc le terme est
$o\left(\frac1{n^\beta}\right)$ : convergente.
\q{$\alpha<1$.} $n\cdot\frac1{n^\alpha\ln n}=\frac{n^{1-\alpha}}{\ln n}\to+\infty$ (croissances comparées) : le terme dépasse $\frac1n$ à partir d'un rang,
divergente.
\q{$\alpha=1$.} $\sum\frac1{n\ln n}$ diverge par comparaison avec une intégrale ($\int_2^X\frac{dt}{t\ln t}=\ln\ln X-\ln\ln2\to+\infty$).
\gain{l'astuce du $\beta$ intermédiaire est la façon standard de « absorber » un logarithme.}
\end{corrige}

\begin{exo}{sn-11}{3}{une série définie par récurrence}
Soit $u_0\in\,]-1,1[$ et $u_{n+1}=\frac{1+u_n^2}2$. On pose $w_n=1-u_n$.
\begin{enumerate}
\item Montrer que $u_n\in[0,1[$ pour $n\geq1$, que $(u_n)$ est croissante, et qu'elle tend vers $1$.
\item Montrer que $w_{n+1}=w_n-\frac{w_n^2}2$. En déduire que $\sum w_n^2$ converge.
\item Déterminer la nature de $\sum\ln\frac{w_{n+1}}{w_n}$, puis celle de $\sum w_n$.
\end{enumerate}
\end{exo}

\begin{corrige}{sn-11}
\q{1.} $u_1=\frac{1+u_0^2}2\in\left[\frac12,1\right[$, et si $u_n\in[0,1[$, alors $u_{n+1}\in\left[\frac12,1\right[$. Ensuite
$u_{n+1}-u_n=\frac{(1-u_n)^2}2\geq0$. Croissante et majorée par $1$, $(u_n)$ converge vers $\ell$ avec $\ell=\frac{1+\ell^2}2$, soit $\ell=1$.
\q{2.} $w_{n+1}=1-\frac{1+(1-w_n)^2}2=w_n-\frac{w_n^2}2$. Donc $w_n^2=2(w_n-w_{n+1})$ : série télescopique, convergente car $(w_n)$ converge
(vers $0$), de somme $2w_0$.
\q{3.} $\ln\frac{w_{n+1}}{w_n}=\ln w_{n+1}-\ln w_n$ et $\ln w_n\to-\infty$ : la série télescopique diverge. Or
$\ln\frac{w_{n+1}}{w_n}=\ln\left(1-\frac{w_n}2\right)\sim-\frac{w_n}2$, de signe constant (négatif) : $\sum w_n$ a la même nature, elle
\textbf{diverge}.
\retenir{une suite $(a_n)$ et la série $\sum(a_{n+1}-a_n)$ ont même nature : c'est un outil pour étudier des suites définies par récurrence.}
\end{corrige}

\begin{exo}{sn-12}{3}{une série trigonométrique surprenante}
\begin{enumerate}
\item Montrer que la série de terme général $u_n=\sin\big((2-\sqrt3)^n\pi\big)$ est à termes positifs et convergente.
\item On pose $A_n=(2-\sqrt3)^n+(2+\sqrt3)^n$. Montrer que $A_n$ est un entier pair (développer par la formule du binôme).
\item En déduire la nature de la série de terme général $v_n=\sin\big((2+\sqrt3)^n\pi\big)$.
\end{enumerate}
\end{exo}

\begin{corrige}{sn-12}
\q{1.} $0<2-\sqrt3<1$, donc $x_n=(2-\sqrt3)^n\pi\in\,]0,\pi[$ pour $n\geq1$ et $u_n>0$. De plus $u_n\sim(2-\sqrt3)^n\pi$, terme d'une série géométrique
convergente : la série converge.
\q{2.} $A_n=\sum_k\binom nk2^{n-k}\big((-\sqrt3)^k+(\sqrt3)^k\big)$ : les termes d'indice $k$ impair s'annulent, ceux d'indice $k=2j$ valent
$2\binom n{2j}2^{n-2j}3^j$. $A_n$ est un entier pair.
\q{3.} $(2+\sqrt3)^n\pi=A_n\pi-(2-\sqrt3)^n\pi$, et $A_n$ est pair : $v_n=\sin\big(-(2-\sqrt3)^n\pi\big)=-u_n$. La série $\sum v_n$
\textbf{converge}, alors que $(2+\sqrt3)^n\to+\infty$.
\end{corrige}

\begin{exo}{sn-13}{4}{reste d'une série alternée}
Soit $R_n=\sum_{k=n+1}^{+\infty}\frac{(-1)^k}k$ le reste de la série $\sum_{k\geq1}\frac{(-1)^k}k$.
\begin{enumerate}
\item Calculer $R_n-R_{n+1}$. Montrer que $R_n+R_{n+1}=\sum_{k=n+1}^{+\infty}\frac{(-1)^k}{k(k+1)}$ et que $|R_n+R_{n+1}|\leq\frac1{(n+1)(n+2)}$.
\item En déduire que $R_n\sim\frac{(-1)^{n+1}}{2n}$, puis que la série $\sum R_n$ converge.
\end{enumerate}
\end{exo}

\begin{corrige}{sn-13}
\q{1.} $R_n-R_{n+1}=\frac{(-1)^{n+1}}{n+1}$. Ensuite $R_{n+1}=\sum_{k=n+1}^{+\infty}\frac{(-1)^{k+1}}{k+1}$ (changement d'indice), donc
$R_n+R_{n+1}=\sum_{k\geq n+1}(-1)^k\left(\frac1k-\frac1{k+1}\right)=\sum_{k\geq n+1}\frac{(-1)^k}{k(k+1)}$. C'est le reste d'une série alternée dont le
terme $\frac1{k(k+1)}$ décroît vers $0$ : il est majoré en valeur absolue par son premier terme $\frac1{(n+1)(n+2)}$.
\q{2.} $2R_n=(R_n-R_{n+1})+(R_n+R_{n+1})=\frac{(-1)^{n+1}}{n+1}+O\left(\frac1{n^2}\right)$, donc $R_n=\frac{(-1)^{n+1}}{2(n+1)}+O\left(\frac1{n^2}\right)$ et
$R_n\sim\frac{(-1)^{n+1}}{2n}$. La série $\sum R_n$ est la somme d'une série alternée convergente et d'une série absolument convergente :
elle converge.
\piege{conclure « $\sum R_n$ converge car $R_n\sim\frac{(-1)^{n+1}}{2n}$ » : l'équivalent ne suffit pas pour une série de signe variable. C'est le
développement $\frac{(-1)^{n+1}}{2(n+1)}+O(\frac1{n^2})$ qui permet de conclure.}
\end{corrige}

\begin{exo}{sn-14}{4}{un produit de Cauchy qui diverge}
Soit $u_n=\frac{(-1)^n}{\sqrt{n+1}}$. Montrer que $\sum u_n$ converge. On note $\sum w_n$ le produit de Cauchy de $\sum u_n$ par elle-même.
Montrer que $|w_n|\geq1$ pour tout $n$ (on pourra majorer $(p+1)(n-p+1)$). Conclure.
\end{exo}

\begin{corrige}{sn-14}
$\sum u_n$ converge par le critère des séries alternées. $w_n=\sum_{p=0}^n\frac{(-1)^p(-1)^{n-p}}{\sqrt{(p+1)(n-p+1)}}=(-1)^n\sum_{p=0}^n\frac1{\sqrt{(p+1)(n-p+1)}}$.
Par l'inégalité $ab\leq\left(\frac{a+b}2\right)^2$, $(p+1)(n-p+1)\leq\left(\frac{n+2}2\right)^2$, donc chacun des $n+1$ termes est
$\geq\frac2{n+2}$ et $|w_n|\geq\frac{2(n+1)}{n+2}\geq1$. Le produit de Cauchy \textbf{diverge grossièrement}.
\retenir{le théorème du produit de Cauchy demande la convergence \emph{absolue} : cet exemple montre qu'on ne peut pas s'en passer.}
\end{corrige}

\begin{exo}{sn-15}{4}{la suite $u_{n+1}=\sin u_n$}
Soit $u_0\in\,]0,\pi[$ et $u_{n+1}=\sin u_n$.
\begin{enumerate}
\item Montrer que $u_n\in\,]0,\pi[$, que $(u_n)$ décroît et tend vers $0$.
\item \emph{Lemme de Cesàro.} Montrer que si $x_n\to\ell\in\R$, alors $\frac1n\sum_{k=1}^nx_k\to\ell$.
\item Montrer que $\frac1{u_{n+1}^2}-\frac1{u_n^2}\to\frac13$. En déduire que $u_n\sim\sqrt{\frac3n}$.
\item Nature de $\sum u_n$, $\sum u_n^2$, $\sum u_n^3$ ?
\end{enumerate}
\end{exo}

\begin{corrige}{sn-15}
\q{1.} Si $u_n\in\,]0,\pi[$, alors $u_{n+1}=\sin u_n\in\,]0,1]\subset\,]0,\pi[$. Comme $\sin x\leq x$ sur $\R_+$, $u_{n+1}\leq u_n$ : la suite décroît,
minorée par $0$, elle converge vers $\ell\in[0,1]$ avec $\sin\ell=\ell$, soit $\ell=0$.
\q{2.} Soit $\varepsilon>0$ et $N$ tel que $|x_k-\ell|\leq\varepsilon$ pour $k>N$. Pour $n>N$ :
$\left|\frac1n\sum_{k=1}^n(x_k-\ell)\right|\leq\frac1n\sum_{k=1}^N|x_k-\ell|+\frac{n-N}n\varepsilon$. Le premier terme tend vers $0$ ($N$ fixé), donc
il est $\leq\varepsilon$ pour $n$ grand : la moyenne est à moins de $2\varepsilon$ de $\ell$.
\q{3.} Avec $x=u_n\to0$ : $\frac1{\sin^2x}-\frac1{x^2}=\frac{x^2-\sin^2x}{x^2\sin^2x}$, et $\sin^2x=x^2-\frac{x^4}3+o(x^4)$, donc le quotient vaut
$\frac{x^4/3+o(x^4)}{x^4+o(x^4)}\to\frac13$. Par Cesàro, $\frac1n\sum_{k=0}^{n-1}\left(\frac1{u_{k+1}^2}-\frac1{u_k^2}\right)=\frac1n\left(\frac1{u_n^2}-\frac1{u_0^2}\right)\to\frac13$,
donc $\frac1{u_n^2}\sim\frac n3$ et $u_n\sim\sqrt{\frac3n}$.
\q{4.} Termes positifs : $u_n\sim\frac{\sqrt3}{n^{1/2}}$ (divergente), $u_n^2\sim\frac3n$ (divergente), $u_n^3\sim\frac{3\sqrt3}{n^{3/2}}$ (convergente).
\end{corrige}
